% Part: incompleteness
% Chapter: representability-in-q
% Section: comp-representable

\documentclass[../../../include/open-logic-section]{subfiles}

\begin{document}

\olfileid{inc}{req}{crq}
\olsection{可计算函数在 $\Th{Q}$ 中可表示}

\begin{thm}
每个可计算函数都在~$\Th{Q}$ 中可表示。
\end{thm}

\begin{proof}
为明确起见，并借助丘奇--图灵论题，我们设定一个函数是可计算的当且仅当它是一般递归的。一般递归函数是指那些能从零函数~$\Zero$、后继函数~$\Succ$ 和投影函数~$\Proj{n}{i}$ 出发，通过复合、原始递归和正则极小化定义的函数。根据 \olref[pri]{lem:prim-rec}，任何使用原始递归从 $f$ 和~$g$ 定义的函数~$h$，也能从 $f$、$g$ 以及 $\Zero$、$\Succ$、$\Proj{n}{i}$、$\Add$、$\Mult$、$\Char{=}$ 出发，通过复合和正则极小化来定义。因此，一个函数是一般递归的，当且仅当它能从 $\Zero$、$\Succ$、$\Proj{n}{i}$、$\Add$、$\Mult$、$\Char{=}$ 出发，通过复合和正则极小化来定义。

此外我们已经证明，上述基本函数在~$\Th{Q}$ 中是可表示的（\cref{inc:req:bre:prop:rep-zero,inc:req:bre:prop:rep-succ,inc:req:bre:prop:rep-proj,inc:req:bre:prop:rep-id,inc:req:bre:prop:rep-add,inc:req:bre:prop:rep-mult}），并且任何由可表示函数通过复合或正则极小化（\olref[cmp]{prop:rep-composition}、\olref[min]{prop:rep-minimization}）定义的函数也是可表示的。因此，每个一般递归函数都在~$\Th{Q}$ 中可表示。
\end{proof}

\begin{explain}
我们已经证明，可计算函数集可以刻画为在 $\Th{Q}$ 中可表示的函数集。事实上，该证明更具一般性。从可表示性的定义不难看出，任何扩张 $\Th{Q}$ 的理论（或者 $\Th{Q}$ 可在其中解释的理论）都能表示可计算函数。但反过来说，在推导关系可计算的任何推导系统中，每个可表示的函数都是可计算的。因此，举例来说，可计算函数集也可以刻画为在皮亚诺算术中、甚至是在策梅洛--Fraenkel 集合论中可表示的函数集。正如 G\"odel 所指出的，这多少有些令人惊讶。我们将会看到，一旦涉及可证性，问题对所考虑的理论非常敏感；粗略地说，公理越强，能证明的命题就越多。然而，在广泛的公理化理论中，可表示的函数恰好就是可计算的函数；只要理论是可公理化的，更强的理论也不会表示更多的函数。
\end{explain}

\end{document}